\subsection{Category-Dependent Classification Response}
\label{sec:mn-categories}

The overall results conceal different category responses. The exported
Line~A class scores contain both gains and losses for every method.
PU-GCN improved 18 categories, left 5 unchanged, and reduced 17; PDANS
improved 16 and reduced 18. Thus a method can lose overall accuracy
without degrading every category. These are descriptive comparisons of
the retained class aggregates, not paired counts of correctly classified
objects.

Table~\ref{tab:mn-category-counts} summarises all 40 categories in each
method--line pair. In Line~B, PU-Net had the most favourable balance:
15 improved, 15 unchanged, and 10 reduced. EAR reduced 24 categories,
while PU-EdgeFormer reduced 20. This broader distribution of losses
helps interpret their lower overall accuracy; it does not identify
which individual objects changed their predictions.

\begin{table}[!t]\centering
\begin{tabular}{llrrrr}\toprule
Line & Method & Improved & Same & Reduced & Mean $\Delta$ (pp)\\\midrule
A & EAR & 16 & 9 & 15 & +0.19 \\
A & PDANS & 16 & 6 & 18 & +0.44 \\
A & PU-Net & 15 & 6 & 19 & -0.54 \\
A & PU-GCN & 18 & 5 & 17 & -0.26 \\
A & PU-EdgeFormer & 17 & 7 & 16 & -1.66 \\
B & EAR & 8 & 8 & 24 & -2.86 \\
B & PDANS & 10 & 12 & 18 & -1.17 \\
B & PU-Net & 15 & 15 & 10 & +0.49 \\
B & PU-GCN & 13 & 8 & 19 & -0.96 \\
B & PU-EdgeFormer & 11 & 9 & 20 & -2.15 \\
\bottomrule\end{tabular}
\caption{Direction of reported category-score changes relative to the
native baseline. Each row covers all 40 categories. The mean is an
unweighted mean of archived category deltas and is not reconstructed OA.}
\label{tab:mn-category-counts}\end{table}

Figures~\ref{fig:mn-cat-A-1}, \ref{fig:mn-cat-A-2}, \ref{fig:mn-cat-B-1}, and~\ref{fig:mn-cat-B-2} retain the complete
category distribution using a common colour range. Categories are split
alphabetically across two panels per line, rather than selected by the
size or sign of a result. For example, Line~B PU-Net gained 25.00~pp on
door but lost 22.62~pp on curtain. Its positive aggregate result therefore
coexists with a substantial category-specific failure. Line~B PU-GCN
gained 11.11~pp on flower pot but lost 16.67~pp on bench. These differences motivate retaining the complete category distribution
in the heatmaps; the underlying scores remain in the archived export.

\thesisfigure{mn_categories_A_1}{Line A category changes, part 1}{Reported class-score changes for Line~A, alphabetical categories 1--20. Blue denotes improvement and red denotes reduction from the native baseline; the colour range is shared by all four maps. PU-EF denotes PU-EdgeFormer. These are archived category aggregates, with no reconstructed object-level counts.}{fig:mn-cat-A-1}
\thesisfigure{mn_categories_A_2}{Line A category changes, part 2}{Reported class-score changes for Line~A, alphabetical categories 21--40. Blue denotes improvement and red denotes reduction from the native baseline; the colour range is shared by all four maps. PU-EF denotes PU-EdgeFormer. These are archived category aggregates, with no reconstructed object-level counts.}{fig:mn-cat-A-2}
\thesisfigure{mn_categories_B_1}{Line B category changes, part 1}{Reported class-score changes for Line~B, alphabetical categories 1--20. Blue denotes improvement and red denotes reduction from the native baseline; the colour range is shared by all four maps. PU-EF denotes PU-EdgeFormer. These are archived category aggregates, with no reconstructed object-level counts.}{fig:mn-cat-B-1}
\thesisfigure{mn_categories_B_2}{Line B category changes, part 2}{Reported class-score changes for Line~B, alphabetical categories 21--40. Blue denotes improvement and red denotes reduction from the native baseline; the colour range is shared by all four maps. PU-EF denotes PU-EdgeFormer. These are archived category aggregates, with no reconstructed object-level counts.}{fig:mn-cat-B-2}

\FloatBarrier

